% -----------------------------------------------------------
\section{Shed}
% -----------------------------------------------------------
	\label{SHED}
	
% % ----------------------
% \subsection{See also}
% % ----------------------

% ----------------------
\subsection{1828 American Dictionary of the English Language} % *1828
% ----------------------
	\label{SHED-1828}

	\begin{compactitem}
		\item[1] To pour out; to effuse; to spill; to suffer to flow out.
		\item[2] To let fall; to cast.
		\item[3] To scatter to emit; to throw off.
	\end{compactitem}

% % ----------------------
% \subsection{Oxford English Dictionary} % *OED
% % ----------------------
	% \label{SHED-OED}

	% \begin{compactitem}
		% \item[]
	% \end{compactitem}

% ----------------------
\subsection{Scriptural Usage} % *SCRIPTURES
% ----------------------
	\label{SHED-SCRIPTURES}

	% % -----
	% \subsubsection{Old Testament} % *OT
	% % -----
		% \label{SHED-OT}
	
		% % --
		% \paragraph{KJV}
		% % --
			% \label{SHED-OT-KJV}
	
			% \begin{tabular}{ll}
				% Total occurrences in the OT & 
			% \end{tabular}
			
			% \begin{compactdesc}
				% \item[]
			% \end{compactdesc}
			
		% % --
		% \paragraph{Hebrew Bible}
		% % --
			% \label{SHED-HEBREW}
		
			% %
			% \subparagraph{\texorpdfstring{\he{sample word}}{}}
			% %
				% \label{PAGA} % whatever is in step bible without periods
			
				% \begin{compactdesc}
					% \item[HALOT , PDF ] \textbf{}
						% \begin{compactitem}
							% \item[1]
						% \end{compactitem}
					% \item[BDB , PDF ] \textbf{}
						% \begin{compactitem}
							% \item[1]
						% \end{compactitem}
					% \item[DCH PDF ] \textbf{}
						% \begin{compactitem}
							% \item[1]
						% \end{compactitem}
				% \end{compactdesc}
				
				% \begin{compactdesc}
					% \item[Some OT uses] \textbf{}
						% \begin{compactdesc}
							% \item[]
						% \end{compactdesc}
				% \end{compactdesc}
			
		% % --
		% \paragraph{Septuagint}
		% % --
			% \label{SHED-LXX}
		
			% %
			% \subparagraph{\texorpdfstring{\gr{sample word}}{}}
			% %
		
	% -----
	\subsubsection{New Testament} % *NT
	% -----
		\label{SHED-NT}
	
		% --
		\paragraph{KJV}
		% --
			\label{SHED-NT-KJV}
			
	
			\begin{tabular}{ll}
				Total occurrences in the NT & 12
			\end{tabular}
			
			\begin{compactdesc}
				\item[Matthew 23:35]
				\item[Matthew 26:26]
				\item[Mark 14:24]
				\item[Luke 11:50] 
				\item[Luke 22:20] 
				\item[Acts 2:33]
				\item[Acts 22:20]
				\item[Romans 3:15]
				\item[{\hyperref[ROMANS-5-5]{Romans 5:5}}]
				\item[Titus 3:6]
				\item[Hebrews 9:22]
				\item[Revelation 16:6]
			\end{compactdesc}
			
		% --
		\paragraph{Greek New Testament} % *GREEK
		% --
			\label{SHED-GREEK}
		
			%
			\subparagraph{\texorpdfstring{\gr{ἐκχέω}}{}}
			%
				\label{EKCHEO}
			
				\begin{compactdesc}
					\item[BDAG 276b] \textbf{}
						\begin{compactitem}
							\item[1] cause to be emitted in quantity, \textit{pour out}
							\item[1a] of liquids
							\item[1b] of solid objects 
							\item[2] cause to fully experience, \textit{pour out} (figurative extension of meaning \#1)...but generally, whatever comes from above as connected with this verb (grace, mercy, etc.)
						\end{compactitem}
					\item[LSJ 526b, PDF 597] \textbf{}
						\begin{compactitem}
							\item[I.1a] pour out
							\item[I.1b] pour away
							\item[I.2] of words, \textit{pour forth, utter}
							\item[I.3] pour out like water, squander, waste
							\item[I.4] spread out
							\item[I.5] throw down
							\item[I.6] shed
							\item[I.7] (Greek synonyms listed)
							\item[II] passive, used by Homer mostly in pluperfect
							\item[II.1] \textit{pour out}, \textit{stream out} or \textit{forth}, properly of liquids...generally, \textit{to be spread out}
							\item[II.2] metaphorical, \textit{to be cast away, forgotten}
							\item[II.3] give oneself up to any emotion, to be overjoyed
							\item[II.4] lie languidly
							\item[II.5] metaphorical, of Time
							\item[II.6] \textit{extend}, of a piece of land
						\end{compactitem}
				\end{compactdesc}
				
				\begin{compactdesc}
					\item[Some NT uses] \textbf{}
						\begin{compactdesc}
							\item[{\hyperref[MATTHEW-26-28]{Matthew 26:28}}] \gr{τοῦτο γάρ ἐστιν τὸ αἷμά μου τῆς διαθήκης τὸ περὶ πολλῶν ἐκχυννόμενον εἰς ἄφεσιν ἁμαρτιῶν·}
							\item[{\hyperref[ROMANS-5-5]{Romans 5:5}}] \gr{ἡ δὲ ἐλπὶς οὐ καταισχύνει· ὅτι ἡ ἀγάπη τοῦ θεοῦ ἐκκέχυται ἐν ταῖς καρδίαις ἡμῶν διὰ πνεύματος ἁγίου τοῦ δοθέντος ἡμῖν.}
								\begin{compactitem}
									\item The use of \gr{ἐκχέω} here is a perfect passive indicative.
								\end{compactitem}
						\end{compactdesc}
				\end{compactdesc}
		
	% -----
	\subsubsection{Book of Mormon} % *BOM
	% -----
		\label{SHED-BOM}
	
		\begin{tabular}{ll}
			Total occurrences in the BOM & 49
		\end{tabular}
		
		\begin{compactdesc}
			\item[1 Nephi 11:22] And I answered him, saying: Yea, it is the love of God, which sheddeth itself abroad in the hearts of the children of men; wherefore, it is the most desirable above all things.
		\end{compactdesc}
		
	% -----
	\subsubsection{Doctrine and Covenants} % *DC
	% -----
		\label{SHED-DC}
	
		\begin{tabular}{ll}
			Total occurrences in the DC & 19
		\end{tabular}
		
		% \begin{compactdesc}
			% \item[]
		% \end{compactdesc}
		
	% -----
	\subsubsection{Pearl of Great Price} % *PGP
	% -----
		\label{SHED-PGP}
	
		\begin{tabular}{ll}
			Total occurrences in the PGP & 3
		\end{tabular}
		
		% \begin{compactdesc}
			% \item[]
		% \end{compactdesc}
		
% ----------------------
\subsection{Key Phrases} % *KEYPHRASES *PHRASES
% ----------------------
	\label{SHED-PHRASES}

	\begin{compactdesc}
		\item[shed abroad] \textbf{2} | Romans 5:5; 1 Nephi 11:22
			\begin{compactdesc}
				\item[Bible anchor verses] Romans 5:5
					\begin{compactdesc}
						\item[Romans 5:5] \gr{ἡ ἀγάπη τοῦ θεοῦ ἐκκέχυται} 
					\end{compactdesc}
				\item[Commentary] See \hyperref[SHED-abroad]{below}.
			\end{compactdesc}
	\end{compactdesc}
	
% % ----------------------
% \subsection{Bible Dictionaries} % *DICTIONARIES
% % ----------------------
	% \label{SHED-BD}

% % ----------------------
% \subsection{Restoration Usage} % *RESTORATION
% % ----------------------
	% \label{SHED-RESTORATION}

	% % -----
	% \subsubsection{Joseph Smith} % *JS
	% % -----
		% \label{SHED-JS}
	
		% \begin{compactdesc}
			% \item[{\hyperref[KEY]{Date/Source}}]
		% \end{compactdesc}
	
	% % -----
	% \subsubsection{Brigham Young} % *BY
	% % -----
		% \label{SHED-BY}
	
		% \begin{compactdesc}
			% \item[{\hyperref[KEY]{Date/Source}}]
		% \end{compactdesc}
	
% ----------------------
\subsection{Commentary and Studies} % *COMMENTARY *STUDIES
% ----------------------
	\label{SHED-COMMENTARY}

	% % -----
	% \subsubsection{Key Points}
	% % -----
		% \label{SHED-KEYS}
	
		% \begin{compactitem}
			% \item
		% \end{compactitem}
		
	% -----
	\subsubsection{``Shed abroad'' or ``shed forth''}
	% -----
		\label{SHED-abroad}
		
		In the KJV, Romans 5:5 reads:
		
		\begin{quote}
			And hope maketh not ashamed; because the love of God is shed abroad in our hearts by the Holy Ghost which is given unto us.
		\end{quote}
		
		Here I'm most interested in the ``the love of God is shed abroad in our hearts'' bit. Here's the Greek:
		
		\begin{quote}
			\gr{ἡ ἀγάπη τοῦ θεοῦ ἐκκέχυται ἐν ταῖς καρδίαις ἡμῶν}
		\end{quote}
		
		The verb is \gr{ἐκκέχυται}, which is a perfect passive form of \gr{ἐκχέω}. In the definitions listed under II, the LSJ says that the word in a passive form has an active kind of intransitive meaning of ``pour out, stream out, stream forth,'' and that the word is often used of liquids. With that in mind, maybe a different translation of that part in Romans 5:5 might be something like this:
		
		\begin{quote}
			The love of God has been streaming forth in(to) our hearts...
		\end{quote}
		
		I'm not sure about the ``in'' vs. ``into'' difference, but I think the key is that the love of God is pouring or streaming into our hearts and then from us outward into the world or into other people. The ``abroad'' in the KJV translation is supposed to capture that, I think, but you might also capture it with the ``in'' translation, where the sense would be more like this:
		
		\begin{quote}
			The love of God has been streaming forth by means of our hearts...
		\end{quote}
		
		The Greek \gr{ἐν} is often used to note the instrument with which something is done, so on this version, it's like our hearts are the means by which that love is both received initially and continues to stream forth into the world and into the lives of other people.
		
		The BOM makes an interesting addition to this idea. Here is the middle part of 1 Nephi 11:22:
		
		\begin{quote}
			Yea, it is the love of God, which sheddeth itself abroad in the hearts of the children of men...
		\end{quote}
		
		The KJV translation of Romans 5:5 is a passive verb with the adverb ``abroad.'' The BOM version, however, is an active verb with a reflexive pronoun, then the adverb. In Spanish a verb with a reflexive pronoun is often used in place of a passive form, but I haven't looked at this enough in the BOM to make a judgment about it. I don't think we need that here, though, because it's more interesting to just take the reflexive pronoun as the object of the verb---the love of God ``sheds itself abroad,'' or \textit{pours itself out} or \textit{streams forth on its own power} in people's hearts. I think the emphasis here is on the love of God's ability to perpetuate its own motion, to multiply itself and duplicate itself from heart to heart. Here, our hearts are more like the vessels that contain the love, but the love of God is the kind of thing that just naturally flows forth under its own power and fills empty spaces. We can participate in that process if we wish, but ultimately it's just going to happen anyway, because of the kind of thing that the love of God is.
		
		See also \hyperref[DC76-53]{DC 76:53}, \hyperref[ACTS-2-33]{Acts 2:33}, 